\DocumentMetadata{
  pdfversion=2.0,
  pdfstandard=ua-2,
  lang=en-US,
  tagging=on,
  tagging-setup={math/setup=mathml-SE,tabsorder=structure}
}
\documentclass[11pt,letterpaper]{article}
\usepackage[margin=0.68in,includefoot]{geometry}
\usepackage{newcomputermodern}
\usepackage{amsmath,amscd,mathtools}
\usepackage{tikz,adjustbox}
\providecommand{\square}{\mdlgwhtsquare}
\usepackage[hidelinks]{hyperref}
\hypersetup{
  pdftitle={Algebra Qualifying Exam, January 2026, Part 1},
  pdfauthor={University of Florida Department of Mathematics},
  pdfsubject={Graduate mathematics examination},
  pdfdisplaydoctitle=true,
  bookmarksopen=true
}
\setcounter{secnumdepth}{0}
\setlength{\parindent}{0pt}
\setlength{\parskip}{0.28em}
\setlength{\emergencystretch}{3em}
\setlength{\leftmargini}{2.15em}
\setlength{\leftmarginii}{2.65em}
\setlength{\leftmarginiii}{3em}
\pagestyle{empty}
\raggedbottom
\allowdisplaybreaks
\NewTaggingSocketPlug{block/list/label}{examproblem}{%
  \tagstructbegin{tag=Lbl}\tagstructend%
  \tagstructbegin{tag=\LBody}%
  \tagstructbegin{tag=\ProblemHeadingTag,title-o={\ProblemHeadingText}}%
  \tagmcbegin{tag=\ProblemHeadingTag,actualtext={\ProblemHeadingText}}%
  #2%
  \tagmcend%
  \tagstructend%
}
\newenvironment{examproblems}[2]{%
  \def\ProblemHeadingTag{#2}%
  \AssignTaggingSocketPlug{block/list/label}{examproblem}%
  \begin{enumerate}%
  \setcounter{enumi}{#1}%
  \renewcommand{\labelenumi}{\textbf{\theenumi.}}%
  \setlength{\topsep}{0.35em}%
  \setlength{\partopsep}{0pt}%
  \setlength{\itemsep}{0.48em}%
  \setlength{\parsep}{0.18em}%
}{\end{enumerate}}
\newenvironment{examsubparts}{%
  \AssignTaggingSocketPlug{block/list/label}{default}%
  \begin{enumerate}%
  \setlength{\topsep}{0.18em}%
  \setlength{\partopsep}{0pt}%
  \setlength{\itemsep}{0.16em}%
  \setlength{\parsep}{0.1em}%
}{\end{enumerate}}
\newenvironment{examinstructions}{%
  \par\begingroup\itshape
}{%
  \par\endgroup\vspace{0.18em}
}
\makeatletter
\renewcommand\subsection{\@startsection{subsection}{2}{0pt}%
  {-1.25ex plus -.25ex minus -.1ex}%
  {0.55ex plus .1ex}%
  {\normalfont\large\bfseries}}
\renewcommand{\@maketitle}{%
  \newpage\null\vskip -1em%
  {\footnotesize\bfseries
    \href{https://www.ufl.edu/}{University of Florida}, \href{https://math.ufl.edu/}{Department of Mathematics}\par}%
  \vskip -0.5em%
  \tagstructbegin{tag=H1,title={Algebra Qualifying Exam, January 2026, Part 1}}%
  \tagpdfparaOff%
  \tagmcbegin{tag=H1}%
    {\LARGE\bfseries\href{https://gma.math.ufl.edu/exams/algebra-qualifying/}{Algebra Qualifying Exam}, January 2026, Part 1\par}%
  \tagmcend%
  \tagpdfparaOn%
  \tagstructend%
  \vskip 0.25em%
  \tagmcbegin{artifact}%
    \hrule height 1pt%
  \tagmcend%
  \vskip 1em%
}
\makeatother
\title{Algebra Qualifying Exam, January 2026, Part 1}
\author{}
\date{}
\begin{document}
\maketitle
\thispagestyle{empty}
\begin{examinstructions}
Answer four problems. Write your answers clearly in complete English sentences. You may quote results (within reason) as long as you state them clearly.
\end{examinstructions}
\begin{examproblems}{0}{H2}
\def\ProblemHeadingText{Problem 1.}
\item
Let \(n \geq 2\) and let~\(G\) be a cyclic group of order~\(n\). Prove that \(\operatorname{Aut}(G) \cong (\mathbb{Z}/n\mathbb{Z})^\times\). (Here \((\mathbb{Z}/n\mathbb{Z})^\times\) is the unit group of the ring \(\mathbb{Z}/n\mathbb{Z}\).)
\def\ProblemHeadingText{Problem 2.}
\item
Prove that there is no simple group of order \(9477=3^6\cdot 13\). (Consider the action of~\(G\) on the set of left cosets of a Sylow~\(3\)-subgroup of~\(G\).)
\def\ProblemHeadingText{Problem 3.}
\item
Construct four groups of order \(28\), and prove that they are pairwise non-isomorphic.
\def\ProblemHeadingText{Problem 4.}
\item
Let~\(G\) be a group.
\begin{examsubparts}
\item[\textnormal{(a)}]
Define the upper central series of~\(G\):
\[
Z_0(G)\leq Z_1(G)\leq Z_2(G)\leq\cdots
\]
\item[\textnormal{(b)}]
Prove that the groups \(Z_i(G)\) are all characteristic subgroups of~\(G\).
\item[\textnormal{(c)}]
Give an example of a nontrivial group~\(G\) such that the groups \(Z_i(G)\) are all trivial.
\end{examsubparts}
\def\ProblemHeadingText{Problem 5.}
\item
Prove that the ring \(\mathbb{Z}[\sqrt{3}]=\{a+b\sqrt{3}:a,b\in\mathbb{Z}\}\) is a Euclidean domain with respect to the norm \(N(a+b\sqrt{3})=|a^2-3b^2|\).
\def\ProblemHeadingText{Problem 6.}
\item
For each of the following polynomials \(f(X)\in\mathbb{Q}[X]\), either prove that \(f(X)\) is irreducible or give a nontrivial factorization of \(f(X)\).
\begin{examsubparts}
\item[\textnormal{(a)}]
\(f(X)=X^4-6X^2+30X-90\)
\item[\textnormal{(b)}]
\(f(X)=X^4+5X^3-6X^2-23\)
\item[\textnormal{(c)}]
\(f(X)=X^4+X^3+X^2+X+1\)
\end{examsubparts}
\end{examproblems}
\end{document}
